{\Large \textbf{Abstract}}

This thesis focuses on the design and implementation of a \verb|Linux| based telemetry
system designed for a custom hydrogen vehicle participating in the \verb|Shell Eco-Marathon| competition.
This system provided a competitive edge over other teams by providing real-time information during the race.

The objective was to create a full hardware-software solution capable of collecting, processing and sending
information to a base station. It uses a \verb|Raspberry Pi| and is able to integrate with peripherals
using \verb|CAN| and \verb|RS485|.

The base system was built using \verb|Yocto| and the telemetry application was written in \verb|C++|.

During the competition the system proved highly useful in troubleshooting various issues. It also provided
the team with data that will help to direct the efforts in future development.

\textbf{Keywords:} Linux, Yocto, C++

\vspace{1cm}

{\Large \textbf{Streszczenie}}

Tematem pracy jest projekt i implementacja systemu telemetrii opartego na systemie Linux,
przeznaczonego dla pojazdu o napędzie wodorowym, który bierze udział w zawodach \verb|Shell Eco-Marathon|.
System ten zapewnił zespołowi przewagę nad konkurencją dzięki dostępowi do danych z pojazdu na żywo w trakcie zawodów.

Celem było stworzenie kompletnego rozwiązania sprzętowo-programowego, zdolnego do gromadzenia, przetwarzania i wysyłania
informacji do stacji bazowej. System korzysta z \verb|Raspberry Pi| i może integrować się z urządzeniami peryferyjnymi,
wykorzystując interfejsy \verb|CAN| i \verb|RS485|.

System bazowy został zbudowany przy pomocy narzędzia \verb|Yocto|, a aplikacja telemetrii została napisana
w języku \verb|C++|.

Podczas zawodów system okazał się niezwykle przydatny w rozwiązywaniu różnych problemów.
Dostarczył równieź danych, które pomogą ukierunkować dalszy rozwój pojazdu.

\textbf{Słowa kluczowe:} Linux, Yocto, C++
